\section{最大共鳴のZariski稠密性}\label{sec:maximal-density}
本節は次数$1$から$6$までを一様に扱う．単純根の共鳴軌跡の稠密性と
有効性の機構はRainsの既知結果\cite[\S4, pp.~34--35]{Filtered}である．
ここでは全根が\emph{同時に}共鳴する点を整数で表示し，
同じ楕円chartでheart一致・不一致の双方が稠密になることを明記する．
数値的な根の共鳴を，全ての根に対する既約sheafの実在と同一視しない．

\subsection{全根共鳴の整数表示}
滑らかな平面三次$E$の変曲点$o$を原点とし，$q\in E$を非torsionとする．
点$q$に対応する次数零線束も$q$と書く．$3\le r\le8$，$d=9-r$とし，
相異なる$p_1,\ldots,p_r\in E$のblow-up $S$を考える．
$h|_E=\OO_E(3o)$，$H=3h-\sum e_i$であり，
\[
 \rho\left(ah-\sum b_ie_i\right)=q^{-m(\alpha)}
 \quad\Longleftrightarrow\quad \sum b_ip_i=m(\alpha)q
\]
を共鳴の符号規約とする．$q$の非torsion性から指数は存在すれば一意である．
根格子の単純根は
\[
 \alpha_0=h-e_1-e_2-e_3,\qquad
 \alpha_i=e_i-e_{i+1}\quad(1\le i<r)
\]
である．

\begin{proposition}[最大共鳴の九価の整数パラメータ]\label{prop:maximal-integral-chart}
全根が$\langle q\rangle$にrestrictionを持つことは，全単純根についてそうであることと同値で，
整数$m_0,\ldots,m_{r-1}$による
\begin{equation}\label{eq:maximal-integral-chart}
 3p_1=(m_0-2m_1-m_2)q,\qquad
 p_i=p_1+\left(\sum_{j=1}^{i-1}m_j\right)q\quad(2\le i\le r)
\end{equation}
で記述される．指数は$m(\alpha_i)=m_i$であり，整数列を固定した式の解は，相異性・genuine性を課す前には
$E[3]$のtorsorで九通りである．相異なる点の範囲では，avatarがgenuineであることと
全根に対する$m(\alpha)\ne0$は同値である．
\end{proposition}
\begin{proof}
単純根が根格子の基底であることから最初の同値を得る．
$e_i-e_{i+1}$のrestrictionは$p_i-p_{i+1}$なので
$p_{i+1}=p_i+m_iq$であり，$p_1+p_2+p_3=m_0q$と合わせれば表示式になる．
複素楕円曲線の$[3]$は全射で核が九個だから，選択の主張も従う．

滑らかな反標準曲線の真の変換も$E$と書く．$H=[E]$はnefで$H^2=d>0$である．
$H$がampleでなければ，Nakai--MoishezonとHodge指数定理から
$HC=0$を満たす既約曲線$C$があり，adjunctionから$C^2=-2$となる．
$C\cap E=\varnothing$なので$\rho(C)=\OO_E$，従って$m(C)=0$である．
逆に$H$がampleなら，任意の根$\alpha$について$H^0(S,\pm\alpha)=0$であり，
$\chi(\alpha-H)=0$とSerre双対性から$H^1(S,\alpha-H)=0$となる．
restriction列より$H^0(E,\rho(\alpha))=0$，従って$\rho(\alpha)\ne\OO_E$である．
非torsion性によりこれは$m(\alpha)\ne0$と同値である．
\end{proof}

\begin{theorem}[全次数の最大共鳴は稠密]\label{thm:maximal-density-main}
相異なる正整数$N_1,\ldots,N_r$に対し$p_i=N_iq$と置く．
この配置のavatarはgenuine del Pezzoで，全ての根は非零整数指数で共鳴する．
さらに，任意の置換$\pi\in\mathfrak S_r$について
\[
 \{(N_1q,\ldots,N_rq):0<N_{\pi(1)}<\cdots<N_{\pi(r)}\}
\]
は$E^r$にZariski稠密である．従って最大共鳴のgenuine locusは稠密であり，
特に$N_1<N_3<N_2$という正負混合条件を課しても稠密である．
\end{theorem}
\begin{proof}
根を$ah-\sum b_ie_i$と書くと
\[
 \sum b_i=3a,\qquad \sum b_i^2=a^2+2.
\]
$a\ge0$の根は，差根$e_i-e_j$，三点を引く直線類，六点を引くconic類，
および$r=8$での$3h-2e_i-\sum_{j\ne i}e_j$に尽きる．
実際Cauchy--Schwarzから$(9-r)a^2\le2r$，
整数$b$に対する$b(b-1)\ge0$から
$\sum b_i(b_i-1)=a^2-3a+2$を得る．
$a=1,2$では全係数が$0,1$，$a=3$では$r=8$で一つが$2$，残りが$1$となる．
唯一残る$a=4,r=8$では，右辺$6$を一つの$3$または$-2$，あるいは
三つの$2$または$-1$で実現しても$\sum b_i=12$に届かず，不可能である．
$a=0$では一つずつの$1,-1$だけである．

従って$m(\alpha)=\sum b_iN_i$は，差根なら相異なる二整数の差，
$a>0$なら正整数の正係数和であり，零にならない．
点の相異性と前命題によりgenuine性を得る．

$D=\{nq:n\ge1\}$は曲線$E$の無限集合なので稠密で，$D^r$も$E^r$に稠密である．
正整数$u_j$を用いて$N_{\pi(j)}=u_1+\cdots+u_j$と書く．
この累積和と置換で定まる$E^r$の群自己同型は整数行列式$\pm1$を持ち，
$D^r$を表示した集合へ送る．よってその集合は稠密である．
$N_1<N_3<N_2$を満たす全順序を一つ選べば最後の主張も従う．
\end{proof}

\begin{center}
\begin{tabular}{c|rrrrrr}
$d$&6&5&4&3&2&1\\\hline
根系&$A_2+A_1$&$A_4$&$D_5$&$E_6$&$E_7$&$E_8$\\
全根数&8&20&40&72&126&240
\end{tabular}
\end{center}
この表は\emph{数値的}共鳴根の数である．同じheartの実在既約根sheafが
この個数あるという意味ではない．整数倍点族では，指定blowdownの$r$個の単純根への
Rainsの存在結果を用い，後述の定理\ref{thm:sorted-spectrum}により全単純とその像を決定する．

\subsection{指定heartの一致条件は構成可能でない}
$U_r\subset E^r$を，相異なる点をblow-upしたavatarがgenuineとなる開集合とする．
固定した$q$と指定順序のRains heartsを各幾何学的点で取る．

\begin{corollary}[一致・不一致の双方の稠密性]\label{cor:nonconstructible-texact}
$U_r$の複素点上で，指定反標準比較がt-exactとなる点集合を$U_r^{\rm ex}$とする．
$U_r^{\rm ex}$と$U_r\setminus U_r^{\rm ex}$はともにZariski稠密で，
いずれも$U_r$の構成可能部分集合の複素点集合ではない．
\end{corollary}
\begin{proof}
定理\ref{thm:sorted-spectrum}は，整数倍点族の指定Rains heartについて，
比較がt-exactであることと$N_1<\cdots<N_r$が同値であることを示す．
昇順の整数配置は前定理により稠密である．
一方$N_1<N_3<N_2$を課した稠密集合では$m(\alpha_2)=N_3-N_2<0$で，
同じ定理，または実在単純根による判定からt-exactでない．
既約noetherian空間の稠密な構成可能部分集合は非空開集合を含む．
一方の集合がそうであれば，他方の稠密性に反する．
\end{proof}

\begin{remark}[稠密性の範囲]
これは固定した楕円chart上の\emph{pointwiseな指定heart}の判定である．
相対的t構造や一様な比較kernelの構成を主張していないので，
そのような相対仮定を備えたt-exactnessの開性定理とは矛盾しない．
また全非可換del Pezzoモジュライの全既約成分に関する稠密性ではない．
共鳴で反標準環のHilbert関数が跳ぶという主張でもない．
$R/gR$が次数$d$の楕円座標環なら，中心元の完全列から依然として
$\dim R_n=1+dn(n+1)/2$である．
\end{remark}
